\section{Mathematical Convention Appendix}
\label{app:mathematical-foundations}

This appendix records only the analytic conventions needed by the continuum
derivations and coordinate appendices. Elementary compactness, Banach-space, and
function-space background is not repeated here.

\begin{convention}[Standing regularity convention]
\label{conv:foundation-standing-regularity}
Unless a section states a stronger or weaker hypothesis, displayed classical
identities are used only under the regularity needed for their terms to be
defined. Fields have the indicated classical derivatives on the stated domains;
boundary traces are assumed when boundary terms are written; and
integration-by-parts or adjoint identities are invoked only when the required
boundary terms are meaningful. When a weak formulation is discussed, the
relevant Sobolev space, trace theorem, and boundary regularity assumptions are
stated locally rather than inferred from this convention. Boundaries are at
least piecewise $C^1$ for classical boundary integrals, and outward unit normals
are interpreted on the regular boundary pieces.
\end{convention}

\begin{definition}[Domain, time, and field notation]
\label{def:foundation-domain-boundary-notation}
For a bounded open connected domain $\Omg\subset\R^n$ with the locally stated
regularity, $\OmgBar$ denotes its closure and $\pOmg$ its boundary. The report
uses
\begin{flalign*}
    \quad
    I_T&\defined (0,T_{\mathrm{final}}),&
    \ITBar&\defined [0,T_{\mathrm{final}}],&
    K&\defined [0,\ell]\times[0,T_{\mathrm{final}}] &&
\end{flalign*}
unless a section declares a local final time or length. Scalar, vector, and
second-order tensor fields are written
\begin{flalign*}
    \quad
    \phi:\Omg\to\R,\qquad
    \vf:\Omg\to\R^n,\qquad
    \vT:\Omg\to\R^{n\times n}. &&
\end{flalign*}
Multi-index notation is used in the standard way:
$D^\alpha\phi$ denotes the classical derivative of order
$|\alpha|=\alpha_1+\cdots+\alpha_n$ when it exists. Spaces such as
$C^k(\OmgBar)$ and $L^p(\Omg)$ are used with their standard meanings; the norm
subscript always identifies the intended space when a norm is needed.
\end{definition}

\begin{definition}[Differential-operator conventions]
\label{def:foundation-differential-operator-conventions}
Named-coordinate partial derivatives use compact notation, for example
\begin{flalign*}
    \quad
    \pd_t\phi=\frac{\pd\phi}{\pd t},\qquad
    \pd_z f=\frac{\pd f}{\pd z},\qquad
    \pd_A p=\frac{\pd p}{\pd A}. &&
\end{flalign*}
For scalar fields, $\nabla\phi$ is the column gradient and
$\del\phi=\nabla\cdot\nabla\phi$. For vector fields, $\nabla\vf$ is the
Jacobian with entries $(\nabla\vf)_{ij}=\pd_{x_j}f_i$ and
\begin{flalign*}
    \quad
    \Div\,\vf=\nabla\cdot\vf=\sum_{i\in[n]}\pd_{x_i}f_i. &&
\end{flalign*}
For tensor fields, the report uses the row-divergence convention
\begin{flalign*}
    \quad
    (\Divbf\vT)_i=\sum_{j\in[n]}\pd_{x_j}T_{ij}. &&
\end{flalign*}
This is the convention used in the continuum momentum balance and in the
Navier--Stokes coordinate appendix. Componentwise integrals and Laplacians are
used for vector- and tensor-valued fields whenever the relevant components are
integrable or differentiable.
\end{definition}

\begin{definition}[Norm conventions]
\label{def:foundation-norm-conventions}
The report uses the usual $L^p$ and $C^k$ norms when they are explicitly named.
For vector-valued $L^p$ quantities, the pointwise range norm is Euclidean. For
matrix or second-order tensor values, the pointwise range norm is the Frobenius
norm,
\begin{flalign*}
    \quad
    \norm{\vT}_{\mathrm F}
    \defined
    \left(\vT:\vT\right)^{1/2}
    =
    \left(\sum_{i,j\in[n]}T_{ij}^{2}\right)^{1/2}. &&
\end{flalign*}
Bare norms are local shorthand only; global notation tables state the intended
space or range norm where ambiguity would otherwise arise.
\end{definition}
